\documentclass[11pt]{article}
\usepackage[a4paper,margin=2.3cm]{geometry}
\usepackage{amsmath,amssymb,bm,amsthm}
\usepackage{booktabs,longtable,array}
\usepackage{enumitem}
\usepackage[dvipsnames]{xcolor}
\usepackage[most]{tcolorbox}
\usepackage{microtype}
\usepackage[colorlinks=true,linkcolor=NavyBlue,citecolor=ForestGreen,urlcolor=NavyBlue]{hyperref}
\setlist{itemsep=2pt,topsep=3pt}
\newtheorem{proposition}{Proposition}
\newtcolorbox{keybox}[1]{colback=NavyBlue!5,colframe=NavyBlue!70,fonttitle=\bfseries,title={#1},left=4pt,right=4pt,top=3pt,bottom=3pt}
\newtcolorbox{warnbox}[1]{colback=BrickRed!5,colframe=BrickRed!70,fonttitle=\bfseries,title={#1},left=4pt,right=4pt,top=3pt,bottom=3pt}
\newtcolorbox{ideabox}[1]{colback=ForestGreen!5,colframe=ForestGreen!70,fonttitle=\bfseries,title={#1},left=4pt,right=4pt,top=3pt,bottom=3pt}
\title{\textbf{Improving User-Centric Clustering for Near-Field Cell-Free XL-MIMO}\\[2pt]
\large Why current rules look the same, which metrics prove a clustering is better, and four proposed strategies}
\author{Prepared for Nayyab Haider (PUC-Rio)}
\date{October 2026}

\begin{document}
\maketitle

\begin{keybox}{Summary in five lines}
\begin{enumerate}[leftmargin=*]
  \item In your code, the channel-norm (CN) and information-rate (IR) clusters are almost identical by construction, so their BER and sum-rate curves overlap. This is a property of the rules, not of the simulation (Section~\ref{sec:audit}).
  \item Clustering rarely changes BER or sum-rate much in your setting, because the cross-AP detector re-reads the non-serving APs whenever the cluster estimate is unreliable. You are right to look for other metrics.
  \item The metrics that \emph{do} separate clustering rules are fronthaul load, gate firing rate, captured information, energy efficiency, tail SE and fairness, and above all the \emph{fronthaul needed to reach a target BER} (Section~\ref{sec:metrics}).
  \item The strongest new strategy for you is a \textbf{detector-aware information-surplus clustering} with a closed-form fronthaul-optimal cluster size. It comes with a provable bound that links the cluster to the information surplus $\Delta_k$ of your own papers (Section~\ref{sec:proposals}).
  \item Show the gain in an \textbf{interference-limited} regime (more users, collinear pairs, sparse APs, estimated CSI), at \textbf{equal fronthaul} or \textbf{equal BER} (Section~\ref{sec:eval}).
\end{enumerate}
\end{keybox}

\begin{warnbox}{Scope of the literature search}
The search used a web search engine. Direct access to arXiv, Google Scholar, IEEE Xplore and the GitHub API was blocked by the environment's network policy, so full texts and code repositories could not be opened. Each reference was checked for title and venue; unverified author lists are marked. Confirm details before citing.
\end{warnbox}

\tableofcontents
\newpage

%=====================================================================
\section{What your code does today (audit)}
\label{sec:audit}

\subsection{Channel-norm (CN) cluster}
For each user $k$ the code forms a metric $m_{l,k}$ and a threshold per propagation regime:
\begin{equation}
  m_{l,k}=\begin{cases}\|\mathbf h^{\rm NF}_{l,k}\|^2\approx N\beta_{l,k}, & A_{l,k}=1,\\ \beta_{l,k}, & A_{l,k}=0,\end{cases}\qquad
  D^{\rm CN}_{l,k}=\mathbb 1\Big[m_{l,k}\ge \tfrac{1}{c}\max_{l'\in\mathcal K^{\rm reg}_k} m_{l',k}\Big],
\end{equation}
with $c=$\,\texttt{clu\_div}$=2$, and the master AP $\arg\max_l m_{l,k}$ is always added.

\textbf{Issue 1 (inconsistent scale).} A near-field metric is about $N\beta$ while a far-field metric is $\beta$. The master AP is therefore almost always a near-field AP, even if a far-field AP has a larger $\beta$. Use the same scale for both regimes (for example $N\beta_{l,k}$, or the estimated $\|\hat{\mathbf h}_{l,k}\|^2$).

\subsection{Information-rate (IR / ``BSR'') cluster}
The code scores each link by
\begin{equation}
  \mathrm{SR}_{l,k}=\log_2\!\Big(1+\frac{p\,(\mathrm{tr}\,\mathbf R_{l,k}-\mathrm{tr}\,\mathbf C_{l,k})}{p\,\mathrm{tr}\,\mathbf C_{l,k}+1}\Big),
\end{equation}
ranks the APs by this score, and keeps \emph{exactly as many APs as the CN cluster has} ($|\mathcal S^{\rm IR}_k|=|\mathcal S^{\rm CN}_k|$).

\textbf{Issue 2 (no interference).} The score has no term for the other users $j\neq k$, so it increases with $\mathrm{tr}\,\mathbf R_{l,k}$ only. With orthogonal pilots and good estimates, $\mathrm{tr}\,\mathbf C\approx0$, and the ranking becomes the ranking by $\beta_{l,k}$, which is the CN ranking. The rate rule of [6], [7] measures the rate \emph{after} the local combiner and includes the interference the combiner leaves; the code version does not.

\textbf{Issue 3 (size tied to CN).} Because the size is copied from the CN cluster and the ranking is (almost) the same, the two clusters coincide. The code's own console line ``AP-selection disagreement: CN vs Rate'' reports this; values near 0 mean every difference between the curves is Monte Carlo noise.

\begin{keybox}{Consequence}
Any new clustering rule you propose must (i) include interference or detector behaviour in its score and (ii) choose its own cluster size. Otherwise it will again coincide with CN.
\end{keybox}

\subsection{Why BER and sum-rate hardly move}
\begin{itemize}
  \item \textbf{The detector compensates.} When the cluster estimate is unreliable, cross-AP list SIC fuses all APs. A weaker cluster makes the gate fire more often, which costs fronthaul, but the BER barely changes. The cost of a bad cluster shows up as \emph{fronthaul}, not as BER.
  \item \textbf{Noise-limited regime.} With $K=8$ users, $N=64$ antennas and $L=8$ APs, the local MMSE combiners already suppress most interference; the earlier GBSE test gave clusters of about 1.3 APs and no visible gain.
  \item \textbf{The sum-rate in Figs.~1/4} is the combiner SINR, which is identical for SIC, list-SIC and cross-AP by construction (they share the combiner).
\end{itemize}

%=====================================================================
\section{Literature review}
\label{sec:lit}

\begin{longtable}{@{}p{0.27\linewidth}p{0.38\linewidth}p{0.28\linewidth}@{}}
\caption{Clustering and AP-selection work relevant to you.}\\
\toprule
\textbf{Work} & \textbf{Idea} & \textbf{Relevance / gap}\\
\midrule
\endfirsthead
\toprule
\textbf{Work} & \textbf{Idea} & \textbf{Relevance / gap}\\
\midrule
\endhead
Bj\"ornson, Sanguinetti [1]; Demir \emph{et al.} [2] & Dynamic cooperation clustering (DCC), master AP and pilot assignment; scalable by construction. & The standard CN baseline; far-field.\\
Buzzi, D'Andrea [5] & User-centric cell-free: each user served by its strongest APs. & Baseline idea.\\
Ngo \emph{et al.} [4] & AP selection by received power to improve total energy efficiency, with a power model including fronthaul. & Source of the EE power model.\\
Dao, Kim [10] & Effective-channel-gain metrics from large-scale fading; sequential user--AP association. & Large-scale only; no detector awareness.\\
Ammar, Adve, Valaee, Yu [15] & Survey of user-centric cell-free: clustering, fronthaul, scalability. & Background and metrics.\\
Mashdour, Salehi, de Lamare, Schmeink, Lima [6]; Mashdour, de Lamare [7] & Clustering based on information rates instead of large-scale fading; fairness scheduling. & Your IR baseline; your code's version omits interference (Issue 2).\\
Tesolin, de Lamare [8] & Graph-based search-and-evaluate (GBSE): serving states as a Hamming graph; discrete search over clusters plus continuous inner evaluation; RAN energy efficiency. & An optimizer you can reuse with a new objective.\\
Tesolin, de Lamare [9] & Graph-based steepest ascent (GBSA) with fractional-programming power allocation; near-optimal RANEE with linear per-iteration cost; compared on processing time, optimality gap and cluster size. & Same; gives the metrics used by your group.\\
G\"ottsch \emph{et al.} [13] & Fixed user-centric clusters under fairness scheduling. & Fairness metrics.\\
HybridUA [14] (authors not verified) & Hybrid network/user-centric association by a z-score of summed large-scale fading; large fronthaul-signalling savings; Jain's index and 5\% outage reported. & Example of fronthaul-first evaluation.\\
G\"ottsch \emph{et al.} [24] & Fronthaul quantization and routing for flexible user load (MILP). & Fronthaul modelling.\\
Ranasinghe \emph{et al.} [11]; Pereira-Ruis\'anchez \emph{et al.} [12] & Graph neural networks for AP selection and cooperation-cluster formation. & Learning-based alternative (future work).\\
Ssettumba \emph{et al.} [16], [17] & Iterative interference cancellation and IDD for clustered cell-free; AP selection with LLR refinement. & Closest to your detector side.\\
Wang \emph{et al.} [19]; Mylonopoulos \emph{et al.} [20] & Near-field cell-free combining and beamfocusing. & No near-field-specific clustering.\\
\bottomrule
\end{longtable}

\textbf{Gap.} Within this search, no clustering rule was found that (a) is designed around a \emph{list/cross-AP detector} and its reliability gate, or (b) uses \emph{near-field range-domain separability} (estimated angle and range) to pick APs. Both are natural for your line of work.

%=====================================================================
\section{Metrics that prove a clustering is better}
\label{sec:metrics}

Below, $\mathcal S_k$ is the cluster of user $k$, $\gamma_{l,k}$ the post-combining SINR of AP $l$ for user $k$ (eq.~(23) in your VTC paper), $\Gamma_k=\sum_{l=1}^L\gamma_{l,k}$ and $\Gamma_{\mathcal S_k}=\sum_{l\in\mathcal S_k}\gamma_{l,k}$.

\begin{enumerate}[label=\textbf{M\arabic*.}]
  \item \textbf{Fronthaul per user} (the most important one for you):
  \begin{equation}
    \bar B=\frac1K\sum_{k=1}^K\big(|\mathcal S_k|+\rho_k\,(L-|\mathcal S_k|)\big),
  \end{equation}
  where $\rho_k$ is the measured gate firing rate (variable \texttt{etaXap\_acc} in your code). A good cluster makes the gate fire less, so it lowers $\bar B$ even when it has the same size.
  \item \textbf{Iso-performance cost}: the fronthaul $\bar B$ (or the energy) needed to reach a target BER, for example $10^{-3}$. This is the cleanest single number: ``at the same BER our clustering needs X\% fewer scalars''.
  \item \textbf{Cluster size} $|\mathcal S_k|$: mean and histogram, separately for near-field-dominated and far-field-dominated users.
  \item \textbf{Captured-information ratio}
  \begin{equation}
    \kappa_k=\frac{\Gamma_{\mathcal S_k}}{\Gamma_k}\in(0,1],
  \end{equation}
  and the information surplus $\Delta_k$ (eq.~(26) of your VTC paper). Plot their CDFs.
  \item \textbf{Gate firing rate} $\rho$ vs.\ SNR (already measured by your code).
  \item \textbf{Tail performance}: 5th-percentile per-user SE (``95\%-likely SE'') and the SE CDF.
  \item \textbf{Fairness}: Jain's index [22]
  \begin{equation}
    J=\frac{\big(\sum_k \mathrm{SE}_k\big)^2}{K\sum_k \mathrm{SE}_k^2}\in\Big[\tfrac1K,1\Big].
  \end{equation}
  \item \textbf{Energy efficiency}, as in [4], [9], with a fronthaul term:
  \begin{equation}
    \mathrm{EE}=\frac{B_w\sum_k \mathrm{SE}_k}{L\,P_{\rm circ}+\sum_l p_l+P_{\rm FH}\,K\bar B},
  \end{equation}
  where $P_{\rm FH}$ is the power per forwarded scalar stream.
  \item \textbf{Complexity}: AP-side combiner cost (per AP, $\mathcal O(N^3)$ times the number of served users if per-user combiners are recomputed), CPU-side LSFD cost $\sum_k\mathcal O(|\mathcal S_k|^3)$, and the run time of the clustering algorithm itself.
  \item \textbf{Load balance}: $\max_l|\mathcal D_l| / \tfrac1L\sum_l|\mathcal D_l|$, where $\mathcal D_l$ is the set of users served by AP $l$.
  \item \textbf{Robustness to CSI error}: the ratio of a metric under estimated vs.\ perfect CSI (your ``retained rate fraction'' figure), now per clustering rule.
\end{enumerate}

\begin{keybox}{Which ones to put in the paper}
M1, M2, M4 (CDF of $\kappa_k$ or $\Delta_k$), M6 and M8 are enough for a strong section. M2 should be the headline figure: BER vs.\ fronthaul at fixed SNR, one curve per clustering rule.
\end{keybox}

%=====================================================================
\section{Proposed clustering strategies}
\label{sec:proposals}

\subsection{P1: Detector-aware information-surplus clustering (main proposal)}
\textbf{Idea.} Choose, for each user, the \emph{smallest} set of APs that captures a fraction $1-\varepsilon$ of the user's total post-combining SINR. Use the estimated channels of \emph{all} AP--UE pairs (your VTC estimator already provides them) and include interference:
\begin{equation}
  \hat\gamma_{l,k}=\frac{p\,|\hat{\mathbf v}_{l,k}^{\mathsf H}\hat{\mathbf h}_{l,k}|^2}
  {p\sum_{j\ne k}|\hat{\mathbf v}_{l,k}^{\mathsf H}\hat{\mathbf h}_{l,j}|^2+\hat{\mathbf v}_{l,k}^{\mathsf H}\big(p\sum_j\mathbf C_{l,j}+\mathbf I\big)\hat{\mathbf v}_{l,k}},
  \qquad
  \mathcal S_k=\arg\min_{\mathcal S}\Big\{|\mathcal S|:\ \sum_{l\in\mathcal S}\hat\gamma_{l,k}\ge(1-\varepsilon)\sum_{l=1}^L\hat\gamma_{l,k}\Big\}.
  \label{eq:p1}
\end{equation}
The minimum is found by sorting $\hat\gamma_{l,k}$ in descending order and adding APs until the target is met ($\mathcal O(L\log L)$ per user).

\begin{proposition}[Cluster target bounds the information surplus]
If $\kappa_k=\Gamma_{\mathcal S_k}/\Gamma_k\ge 1-\varepsilon$, then the information surplus of your papers satisfies
\begin{equation}
  \Delta_k=\frac{\Gamma_k-\Gamma_{\mathcal S_k}}{(1+\Gamma_{\mathcal S_k})(1+\Gamma_k)}\;\le\;\frac{\varepsilon}{1+(1-\varepsilon)\Gamma_k}.
\end{equation}
\end{proposition}
\begin{proof}
$\Gamma_k-\Gamma_{\mathcal S_k}\le\varepsilon\Gamma_k$ and $1+\Gamma_{\mathcal S_k}\ge1+(1-\varepsilon)\Gamma_k$, so $\Delta_k\le \varepsilon\Gamma_k/\big((1+(1-\varepsilon)\Gamma_k)(1+\Gamma_k)\big)$, and $\Gamma_k/(1+\Gamma_k)\le1$.
\end{proof}
\textbf{Why this is a good contribution.} One parameter $\varepsilon$ directly controls how much the cluster-only estimate can lose against the all-AP estimate. It links clustering to the detector analysis of your ICASSP and VTC papers. The cluster size adapts per user: a near-field user whose energy sits in one or two APs gets a small cluster, while an edge user gets a larger one. This is where fronthaul savings come from.

\subsection{P2: Fronthaul-optimal cluster size for the gated cross-AP detector}
\textbf{Idea.} Your detector forwards the $|\mathcal S_k|$ cluster scalars always, and the remaining $L-|\mathcal S_k|$ only when the gate fires. A larger cluster costs more base fronthaul but makes the gate fire less. There is an optimum.

\textbf{Gate model.} The code declares a decision unreliable when the soft estimate is farther than $d_{\rm abs}=d_{\rm th}\,d_{\min}/2$ from every constellation point. If the fused estimate error is $\mathcal{CN}(0,\varepsilon^2_{\mathcal S})$ with $\varepsilon^2_{\mathcal S}=1/(1+\Gamma_{\mathcal S})$ (eq.~(25) of VTC), the probability of staying inside the disk around the true point is $1-e^{-d_{\rm abs}^2/\varepsilon_{\mathcal S}^2}$, so
\begin{equation}
  \rho_k(\mathcal S)\approx\exp\!\big(-d_{\rm abs}^2\,(1+\Gamma_{\mathcal S})\big).
\end{equation}
\textbf{Optimal size.} With the APs sorted by $\hat\gamma_{l,k}$ and $\mathcal S^{(m)}$ the $m$ best,
\begin{equation}
  m_k^\star=\arg\min_{m\in\{1,\dots,L\}}\Big[m+\exp\!\big(-d_{\rm abs}^2(1+\Gamma_{\mathcal S^{(m)}})\big)(L-m)\Big],
  \label{eq:p2}
\end{equation}
a one-dimensional search of $L$ values per user. P1 and P2 can be combined: P2 picks the size, P1's ranking picks the APs.

\textbf{Why this is a good contribution.} It is closed-form, cheap, and specific to the gated cross-AP architecture you proposed. It also predicts your fronthaul figure, so you can validate it: compare the model $\rho_k(\mathcal S)$ with the measured \texttt{etaXap\_acc}.

\subsection{P3: Near-field separability-aware AP ranking}
\textbf{Idea.} In the near field, an AP is valuable for user $k$ if it can \emph{separate} $k$ from its strongest interferers in angle \emph{or range}. The polar-domain estimator gives estimated angles and ranges $(\hat\theta_{l,j},\hat r_{l,j})$ for every pair, so the CPU can predict the near-field correlation
\begin{equation}
  \hat\mu_l(k,j)=\big|\mathbf b(\hat\theta_{l,k},\hat r_{l,k})^{\mathsf H}\mathbf b(\hat\theta_{l,j},\hat r_{l,j})\big|
\end{equation}
and score
\begin{equation}
  \mathrm{score}_{l,k}=\hat\gamma_{l,k}\Big(1-\max_{j\in\mathcal I_k}\hat\mu_l(k,j)^2\Big),
\end{equation}
where $\mathcal I_k$ are the strongest co-channel users of $k$. Your simulations deliberately create \emph{collinear user pairs}, where this score should help most. With sparse APs (see the sparse-array document), $\hat\mu$ also captures grating-lobe ambiguities.

\subsection{P4: Graph-based search with the new objective (link to GBSE/GBSA)}
Use the GBSE/GBSA framework of [8], [9] as the optimizer, but with an objective that includes the detector:
\begin{equation}
  J(\{\mathcal S_k\})=\sum_k \mathrm{SE}_k-\lambda\,K\bar B\quad\text{or}\quad\mathrm{EE}\ \text{(M8) with the gate term}.
\end{equation}
Start the search from the P1/P2 solution (warm start), which cuts the number of iterations. This shows continuity with your group's work and gives a near-optimal benchmark for P1/P2.

\subsection{Recommendation}
\begin{keybox}{What to propose}
Make \textbf{P1 + P2} the proposed method (analytical, $\mathcal O(L\log L)$ per user, with Proposition~1 and the fronthaul model). Add \textbf{P3} as the near-field enhancement. Use \textbf{P4} (GBSE with the new objective) as an upper-performance benchmark. A GNN version can be future work.
\end{keybox}

\begin{keybox}{Algorithm 1 --- P1+P2: detector-aware clustering (per coherence block, at the CPU)}
\textbf{Input:} estimates $\hat{\mathbf h}_{l,k}$, error covariances $\mathbf C_{l,k}$, power $p$, target $\varepsilon$, gate radius $d_{\rm abs}$.
\begin{enumerate}[leftmargin=*]
  \item For every AP $l$: compute the local MMSE combiners $\hat{\mathbf v}_{l,k}$ for all users and the SINRs $\hat\gamma_{l,k}$ from \eqref{eq:p1}.
  \item For every user $k$:
  \begin{enumerate}
    \item sort the APs so that $\hat\gamma_{l_1,k}\ge\hat\gamma_{l_2,k}\ge\dots$, and form the cumulative sums $\Gamma^{(m)}=\sum_{i\le m}\hat\gamma_{l_i,k}$;
    \item P1: $m_\varepsilon=\min\{m:\Gamma^{(m)}\ge(1-\varepsilon)\Gamma^{(L)}\}$;
    \item P2: $m^\star=\arg\min_m\big[m+e^{-d_{\rm abs}^2(1+\Gamma^{(m)})}(L-m)\big]$;
    \item $\mathcal S_k=\{l_1,\dots,l_{\max(m_\varepsilon,m^\star)}\}$ (use $\min$ instead of $\max$ for a fronthaul-first design).
  \end{enumerate}
\end{enumerate}
\textbf{Cost:} $\mathcal O(L\log L)$ per user after the combiners, which the receiver computes anyway.
\end{keybox}

%=====================================================================
\section{How to show the gain (evaluation design)}
\label{sec:eval}

\paragraph{Baselines.} Single-AP (lower bound), CN/DCC [1], IR [6], effective-gain AP selection [10], GBSE/GBSA with RANEE [8], [9], full cooperation (upper bound for BER, worst for fronthaul).

\paragraph{Regime.} Clustering only matters when interference matters. Use at least one of: more users ($K=16$--$32$); fewer antennas per AP ($N=16$--$32$); the collinear pairs you already generate; sparse APs; estimated CSI. Report the regime honestly.

\paragraph{Figures.}
\begin{enumerate}
  \item \textbf{BER vs.\ fronthaul} $\bar B$ at fixed SNR (sweep $\varepsilon$ for P1, the threshold for CN, the size for IR): the headline trade-off.
  \item $\bar B$ and gate firing rate $\rho$ vs.\ SNR, per rule; model of P2 vs.\ measured $\rho$.
  \item CDF of $\kappa_k$ and of $\Delta_k$ per rule (validates Proposition~1).
  \item Per-user SE CDF and Jain's index.
  \item Cluster-size histogram, near-field vs.\ far-field users.
  \item EE vs.\ maximum cluster size (as in GBSA [9]).
  \item Run time of the clustering algorithms vs.\ $K$.
\end{enumerate}

%=====================================================================
\section{Implementation notes for your MATLAB code}
\label{sec:impl}
\begin{enumerate}
  \item Compute the clusters \emph{after} channel estimation in each SNR loop, from $\hat{\mathbf h}$ and $\mathbf C$ (as \texttt{D\_BSR\_s} is computed now). Never use \texttt{h\_det\_all} in a clustering rule for an estimated-CSI result; with \texttt{genie\_csi = false} the estimates are authentic.
  \item Do \emph{not} tie the cluster size to \texttt{D\_CN}; let each rule choose its size.
  \item Fix the master-AP scale (Issue 1).
  \item Reuse what the code already measures: \texttt{etaXap\_acc} ($\rho$), \texttt{loadBSR\_acc} (size), \texttt{mBERrt*} (per-user BER for SE), and add $\kappa_k$, $\Delta_k$ and Jain's index.
  \item Run the clustering rules on the \emph{same} channel and noise realizations, so that differences are not Monte Carlo noise.
\end{enumerate}

\begin{warnbox}{Honest risks}
\begin{itemize}
  \item In the default regime ($K=8$, $N=64$) all rules may give the same BER. That is expected; the gain must be shown in fronthaul and EE at equal BER.
  \item The gate model in P2 assumes a Gaussian fused error and ignores the disks around other constellation points; it is an approximation. Validate it against the measured rate before relying on it.
  \item P3 depends on the accuracy of the estimated angle and range; report its sensitivity to estimation error.
\end{itemize}
\end{warnbox}

%=====================================================================
\begin{thebibliography}{99}\small
\bibitem{bjsc} E.~Bj\"ornson and L.~Sanguinetti, ``Scalable cell-free massive MIMO systems,'' \emph{IEEE Trans.\ Commun.}, vol.~68, no.~7, pp.~4247--4261, 2020.
\bibitem{found} \"O.~T.~Demir, E.~Bj\"ornson, and L.~Sanguinetti, ``Foundations of user-centric cell-free massive MIMO,'' \emph{Found.\ Trends Signal Process.}, vol.~14, no.~3--4, pp.~162--472, 2021.
\bibitem{ngo17} H.~Q.~Ngo, A.~Ashikhmin, H.~Yang, E.~G.~Larsson, and T.~L.~Marzetta, ``Cell-free massive MIMO versus small cells,'' \emph{IEEE Trans.\ Wireless Commun.}, vol.~16, no.~3, pp.~1834--1850, 2017.
\bibitem{ngo18} H.~Q.~Ngo, L.-N.~Tran, T.~Q.~Duong, M.~Matthaiou, and E.~G.~Larsson, ``On the total energy efficiency of cell-free massive MIMO,'' \emph{IEEE Trans.\ Green Commun.\ Netw.}, vol.~2, no.~1, pp.~25--39, 2018.
\bibitem{buzzi} S.~Buzzi and C.~D'Andrea, ``Cell-free massive MIMO: User-centric approach,'' \emph{IEEE Wireless Commun.\ Lett.}, vol.~6, no.~6, pp.~706--709, 2017.
\bibitem{mash} S.~Mashdour, S.~Salehi, R.~C.~de~Lamare, A.~Schmeink, and J.~P.~S.~H.~Lima, ``Clustering and scheduling with fairness based on information rates for cell-free MIMO networks,'' \emph{IEEE Wireless Commun.\ Lett.}, vol.~13, no.~7, pp.~1798--1802, 2024.
\bibitem{mash2} S.~Mashdour and R.~C.~de~Lamare, ``Study of clustering techniques and scheduling algorithms with fairness for cell-free MIMO networks,'' arXiv:2404.18032, 2024.
\bibitem{gbse} J.~C.~C.~Tesolin and R.~C.~de~Lamare, ``Mixed-integer optimization based on graph-based decomposition for user-centric cell-free networks,'' \emph{IEEE Wireless Commun.\ Lett.}, vol.~15, 2026, doi:10.1109/LWC.2026.3700959.
\bibitem{gbsa} J.~C.~C.~Tesolin and R.~C.~de~Lamare, ``Study of graph-based search for energy-efficient clustering in cell-free massive MIMO networks,'' arXiv:2607.04074, 2026.
\bibitem{dao} H.~T.~Dao and S.~Kim, ``Effective channel gain-based access point selection in cell-free massive MIMO systems,'' \emph{IEEE Access}, vol.~8, pp.~108127--108132, 2020.
\bibitem{gnn1} V.~Ranasinghe, N.~Rajatheva, and M.~Latva-aho, ``Graph neural network based access point selection for cell-free massive MIMO systems,'' in \emph{Proc.\ IEEE GLOBECOM}, 2021; arXiv:2107.02884.
\bibitem{gnn2} Pereira-Ruis\'anchez \emph{et al.}, ``A GNN-based approach to AP cooperation cluster formation in cell-free massive MIMO,'' in \emph{Proc.\ IEEE VTC2025-Spring}, doi:10.1109/VTC2025-SPRING65109.2025.11174418 (initials not verified).
\bibitem{gott1} F.~G\"ottsch \emph{et al.}, ``User-centric clustering under fairness scheduling in cell-free massive MIMO,'' arXiv:2305.08363, 2023 (full author list not verified).
\bibitem{hybridua} ``Hybrid network- and user-centric scalable cell-free massive MIMO for fronthaul signaling minimization,'' arXiv:2412.15475, 2024 (authors not verified).
\bibitem{ammar} M.~Ammar, R.~Adve, S.~Valaee, and W.~Yu, ``User-centric cell-free massive MIMO networks: A survey of opportunities, challenges and solutions,'' \emph{IEEE Commun.\ Surveys Tuts.}, vol.~24, no.~1, pp.~611--652, 2022.
\bibitem{sset1} T.~Ssettumba, S.~Mashdour, L.~T.~N.~Landau, P.~B.~da~Silva, and R.~C.~de~Lamare, ``Iterative interference cancellation for clustered cell-free massive MIMO networks,'' \emph{IEEE Wireless Commun.\ Lett.}, vol.~14, no.~2, pp.~509--513, 2024.
\bibitem{sset2} T.~Ssettumba, Z.~Shao, L.~T.~N.~Landau, M.~S.~P.~Facina, P.~R.~B.~da~Silva, and R.~C.~de~Lamare, ``Centralized and decentralized IDD schemes for cell-free massive MIMO systems: AP selection and LLR refinement,'' \emph{IEEE Access}, vol.~12, pp.~62392--62406, 2024.
\bibitem{own} N.~Haider and R.~C.~de~Lamare, ``Cross-AP list-based interference cancellation for near-field cell-free XL-MIMO systems'' (ICASSP submission) and ``Cross-AP channel estimation and interference cancellation for near-field cell-free XL-MIMO systems'' (VTC submission), 2026.
\bibitem{wangcomb} Z.~Wang, J.~Zhang, B.~Xu, D.~Niyato, B.~Ai, S.~Mao, and Z.~Han, ``Low-complexity distributed combining design for near-field cell-free XL-MIMO systems,'' \emph{IEEE Trans.\ Wireless Commun.}, vol.~25, pp.~11799--11815, 2026.
\bibitem{mylo} G.~Mylonopoulos, G.~Interdonato, S.~Buzzi, and P.~Liu, ``Leveraging line-of-sight propagation for near-field beamfocusing in cell-free networks,'' arXiv:2503.21599, 2025.
\bibitem{cui} M.~Cui and L.~Dai, ``Channel estimation for extremely large-scale MIMO: Far-field or near-field?'' \emph{IEEE Trans.\ Commun.}, vol.~70, no.~4, pp.~2663--2677, 2022.
\bibitem{jain} R.~Jain, D.-M.~Chiu, and W.~Hawe, ``A quantitative measure of fairness and discrimination for resource allocation in shared computer systems,'' DEC Tech.\ Rep.\ TR-301, 1984.
\bibitem{haiderwcl} N.~Haider and R.~C.~de~Lamare, ``Iterative interference cancellation for mixed near- and far-field XL-MIMO systems,'' \emph{IEEE Wireless Commun.\ Lett.}, vol.~15, pp.~5323--5327, 2026.
\bibitem{gott2} F.~G\"ottsch \emph{et al.}, ``Optimizing fronthaul quantization for flexible user load in cell-free massive MIMO,'' arXiv:2510.06734, 2025 (full author list not verified).
\end{thebibliography}

\end{document}
